\section{Property Composition and Concept Discovery}
\label{sec:composition}

Sections~\ref{sec:property}--\ref{sec:cross} analyzed properties individually. This section asks: can the system \emph{construct} new properties from existing ones, and can those constructions reveal concepts that no individual property captures?

\subsection{Conjunction}

$\texttt{define\_compose\_and}(\textit{name}, p, q)$ constructs a new property whose extension is $\mathrm{ext}(p) \cap \mathrm{ext}(q)$. Three rules are generated:

\begin{align*}
\mathit{has\_property}(x, p) \wedge \mathit{has\_property}(x, q) &\vdash \mathit{has\_property}(x, \textit{name}) & \text{(detect)} \\
\mathit{has\_property}(x, \textit{name}) &\vdash \mathit{has\_property}(x, p) & \text{(split-left)} \\
\mathit{has\_property}(x, \textit{name}) &\vdash \mathit{has\_property}(x, q) & \text{(split-right)}
\end{align*}

On $\mathbb{Z}_3$, $\textit{two\_sided\_id} = \textit{left\_id} \wedge \textit{right\_id}$ has extension $\{e_0\}$, identical to both parents. Phase~6 correctly discovers $\textit{two\_sided\_id} \leftrightarrow \textit{left\_id} \leftrightarrow \textit{right\_id}$.

On the magma (no right identity), the conjunction has empty extension, and no false equivalences are derived.

\subsection{Negation}

$\texttt{define\_negate}(\textit{name}, p)$ constructs a property whose extension is the complement: $\mathrm{ext}(\textit{name}) = \mathrm{domain} \setminus \mathrm{ext}(p)$. Detection uses stratified negation at stratum~3:

\[
\mathit{element}(x),\; \lnot\, \mathit{has\_property}(x, p) \;\vdash\; \mathit{has\_property}(x, \textit{name})
\]

Negation is the key to escaping the ``0 informative combinations'' result of pure conjunction. Without negation, all property pairs on the test domains produce degenerate combinations (extension equals one parent, or is empty). With negation, orthogonal properties emerge.

\subsection{Combination Scoring}

Not all conjunctions are informative. A scoring function filters degenerate combinations:

\paragraph{Hard filters (score = 0):}
\begin{itemize}
\item \textbf{Empty extension}: $\mathrm{ext}(p \wedge q) = \varnothing$.
\item \textbf{Degenerate to parent}: $\mathrm{ext}(p \wedge q) = \mathrm{ext}(p)$ or $\mathrm{ext}(q)$.
\item \textbf{Cross-degenerate}: $\mathrm{ext}(p \wedge q) = \mathrm{ext}(r)$ for some existing property $r$.
\item \textbf{Trivial universal}: $\mathrm{ext}(p \wedge q) = \mathrm{domain}$.
\end{itemize}

\paragraph{Soft scoring (score $> 0$):}
\[
\text{score} = \underbrace{\min(1 - J_p,\; 1 - J_q)}_{\text{independence}} \;\times\; \underbrace{4r(1-r)}_{\text{rarity}} \;\times\; \underbrace{1 - J(p, q)}_{\text{constraint}}
\]
where $J_p = J(\mathrm{ext}(p \wedge q),\, \mathrm{ext}(p))$ is the Jaccard similarity between the combination and each parent, and $r = |\mathrm{ext}(p \wedge q)| / |\mathrm{domain}|$.

\subsection{Results: Autonomous Concept Discovery}

\paragraph{Without negation (pure conjunction).} Across all three domains (algebra, order theory, set theory), auto-enumeration of all property pairs finds \textbf{0 informative combinations}. Every conjunction is either empty, degenerates to a parent, or equals an existing property. This is a genuine negative result: the property sets are already ``minimally complete'' for these structures.

\paragraph{With negation.} On a 4-element chain $a < b < c < d$:

\begin{table}[h]
\centering
\caption{Auto-discovered combinations on a 4-element chain, after negation. Cross-degenerate filter applied.}
\label{tab:combos}
\begin{tabular}{llcl}
\toprule
Combination & Extension & Score & Interpretation \\
\midrule
$\lnot\text{maximal} \wedge \lnot\text{minimal}$ & $\{b, c\}$ & 0.167 & interior elements \\
\midrule
\multicolumn{4}{l}{\emph{Filtered by cross-degenerate check:}} \\
$\text{extremal} \wedge \lnot\text{minimal}$ & $\{d\} = \mathrm{ext}(\text{is\_maximal})$ & 0 & redundant \\
$\lnot\text{maximal} \wedge \text{extremal}$ & $\{a\} = \mathrm{ext}(\text{is\_minimal})$ & 0 & redundant \\
\bottomrule
\end{tabular}
\end{table}

The sole surviving combination, $\lnot\text{maximal} \wedge \lnot\text{minimal}$, identifies the \emph{interior elements} of the chain---a recognized mathematical concept that the system discovers without being told its name or definition. The cross-degenerate filter correctly eliminates two combinations that merely reconstruct existing properties (\text{is\_maximal} and \text{is\_minimal}) through negation indirection.

\subsection{Interpretation}

\begin{enumerate}
\item \textbf{Negation as orthogonality generator.} Pure conjunction over a ``complete'' property set yields nothing new. Negation creates properties that are by construction orthogonal to existing ones, opening the combinatorial space.

\item \textbf{Cross-degenerate filtering is essential.} Without it, 2 of 3 candidates were redundant reconstructions of existing properties. The filter reduces false discovery rate from 67\% to 0\%.

\item \textbf{Carrier size matters.} On 3-element carriers, even negation cannot escape degeneracy (all property subsets form chains). The 4-element chain provides just enough room for a genuinely new intersection. This suggests that concept discovery requires a minimum structural complexity in the input domain.
\end{enumerate}

\subsection{Disjunction}

$\texttt{define\_compose\_or}(\textit{name}, p, q)$ constructs a property with $\mathrm{ext}(\textit{name}) = \mathrm{ext}(p) \cup \mathrm{ext}(q)$. Only forward rules are generated (no negation-based decomposition):
\begin{align*}
\mathit{has\_property}(x, p) &\vdash \mathit{has\_property}(x, \textit{name}) \\
\mathit{has\_property}(x, q) &\vdash \mathit{has\_property}(x, \textit{name})
\end{align*}

Scoring follows the same filter structure as conjunction, with soft factors adapted for unions: $\text{score} = \text{covering\_penalty} \times \text{disjointness} \times \text{rarity}$.

On the 4-element chain, $\mathit{or}(\text{maximal}, \text{minimal}) = \{a, d\} = \mathrm{ext}(\text{extremal})$, correctly filtered as \texttt{degenerate\_to\_is\_extremal}.

% Property space analysis and meta-level observations are in Section \ref{sec:space}.
